\section{Results}
\label{sec:results}

Everything below is the untouched test window, 2018-07-01 to 2018-08-31, scored
once for models already chosen. Cross-validated and validation columns sit beside
it because the disagreement between them is part of the result. Figures are in
Appendix~\ref{app:figures}.

\subsection{Problem 1A --- lead time}
\label{sec:results-reg}

\begin{table}[htbp]
\centering
\caption{Lead-time regression: two reference points that are not models, then
each family's complexity ladder. \best{Bold} is the selected model.
Figure~\ref{fig:ladder}.}
\label{tab:reg}
\small
\begin{tabular}{@{}llrrrrr@{}}
\toprule
\textbf{Model} & \textbf{Family} & \textbf{CV MAE} & \textbf{Val MAE} &
\textbf{Test MAE} & \textbf{RMSE} & \textbf{$R^2$} \\
\midrule
Mean of the training period      & baseline & ---  & ---  & 6.645 & 7.572 & $-0.990$ \\
Olist's own delivery promise     & baseline & ---  & ---  & 9.788 & 12.155 & $-4.130$ \\
\midrule
LinearRegression \emph{(rung 0)} & A linear & 5.242 & 4.965 & 4.181 & 5.639 & $-0.104$ \\
Ridge \emph{(rung 1)}            & A linear & 5.201 & 5.164 & 4.309 & 5.775 & $-0.158$ \\
DecisionTree \emph{(rung 0)}     & B tree   & 5.539 & 5.927 & 6.116 & 8.664 & $-1.606$ \\
RandomForest \emph{(rung 1)}     & B tree   & 5.457 & 5.513 & 4.956 & 6.169 & $-0.321$ \\
HistGradientBoosting \emph{(rung 2)} & B tree & 5.495 & 4.802 & 4.141 & 5.517 & $-0.057$ \\
\best{VotingRegressor (A+B)}     & C ensemble & 5.315 & \best{4.784} & \best{3.898} & \best{5.342} & \best{$+0.009$} \\
StackingRegressor (A+B)          & C ensemble & 5.443 & 4.792 & 4.121 & 5.430 & $-0.024$ \\
\bottomrule
\end{tabular}
\end{table}

The selected model puts mean absolute error at \textbf{3.90 days}: \textbf{41.3\%
below} a constant fitted on the training period and \textbf{60.2\% below the
platform's own promised date}. That second comparison is the one operations cares
about, and its direction is not anomalous --- a promised date is a commercial
decision, padded because a broken promise costs returns and retention
\citep{cui2024promise}. It is an excellent feature and a poor forecast.

Three things matter more than the winner. \textbf{Regularisation does not pay:}
Ridge is \emph{worse} than plain least squares on validation (5.164 against
4.965) and on test (4.309 against 4.181), with a tuned $\alpha$ of 660 that
cross-validation preferred and no later window did --- within family A, rung 1
fails to beat rung 0. \textbf{Inside family B complexity does pay}, monotonically
on validation: 5.927 $\rightarrow$ 5.513 $\rightarrow$ 4.802. And \textbf{only
the ensemble has a positive $R^2$}; every other model is negative, which is the
signature of a level shift rather than of a useless model, and the reason MAE
leads throughout. The blend beats both its members (4.181 and 4.141 against
3.898) by a specific mechanism: a tree cannot predict outside the range of the
targets it was fitted on, because every leaf averages training values, and
test-window lead times are systematically shorter than anything common in
training. A linear model has no such bound and extrapolates the level
(Figure~\ref{fig:extrap}). Averaging is one member covering the other's exact
failure, which is why the weights are fixed at 1:2 in the tree's favour rather
than tuned.

\subsection{Problem 1B --- late against the promise}
\label{sec:results-late}

\begin{table}[htbp]
\centering
\caption{The late-order flag, best of each route; the full seven-candidate table
is in \code{report/results/clf\_p1b.csv}. No-skill PR-AUC floor 0.074.
Figure~\ref{fig:routes}.}
\label{tab:late}
\small
\begin{tabular}{@{}llrrrrr@{}}
\toprule
\textbf{Route} & \textbf{Model} & \textbf{ROC-AUC} & \textbf{PR-AUC} &
\textbf{lift} & \textbf{bal.\ acc} & \textbf{F1} \\
\midrule
direct  & best of five (HistGradientBoosting) & 0.584 & 0.092 & 1.24$\times$ & 0.524 & 0.134 \\
derived & predicted days $>$ promised days, unfitted & 0.569 & 0.109 & 1.47$\times$ & 0.569 & 0.214 \\
\best{derived} & \best{regression margin, calibrated} & \best{0.782} & \best{0.214} & \best{2.89$\times$} & \best{0.705} & \best{0.235} \\
\bottomrule
\end{tabular}
\end{table}

The deliverable is the derived route, and that follows from the framing fixed in
Section~\ref{sec:problem} rather than from the numbers: Problem 1 is the brief's
dual framing, so the flag \emph{is} the regression's own answer compared against
the promised date. The five direct classifiers answer a different question ---
does the flag need a model of its own? It does not. The derived route reaches
\textbf{2.89$\times$} the no-skill floor where the best direct classifier reaches
1.24$\times$, and it costs no additional model: the margin
$\hat{y} - \text{promised}$ is mapped to a probability by a single-parameter
logistic \citep{platt1999probabilistic} fitted on validation. Even the
zero-parameter rule --- flag whenever the model expects a miss --- beats every
fitted classifier.

\paragraph{The selection protocol got this one wrong, and that is worth reporting.}
Applied to these seven candidates the rolling-origin rule of Section~\ref{sec:cv}
ranks a direct classifier first. On every in-training window the candidates sit
within noise of one another --- in the fifth fold they span 0.206 to 0.362, with
the derived route mid-pack --- so the rule has almost nothing to separate them
on, and separates them on noise; on the forward window the gap opens to a factor
of 2.3. The monthly series explains why: the later training folds and the
validation period are dominated by the February--April 2018 congestion, when the
late rate ran at 16.0\% and 21.4\%, and by the May strike. A classifier can fit a
regime; the derived route encodes a mechanism, and mechanisms transfer. We report
the disagreement rather than presenting the rule as having chosen this, and note
that confirming the preference would need a further forward window we do not
have.

\subsection{Problem 2 --- low-review risk}
\label{sec:results-review}

\begin{table}[htbp]
\centering
\caption{Low-review classification at delivery. Accuracy is printed beside the
majority-class rule that scores 0.890 while catching nothing. No-skill PR-AUC
floor 0.110. Figure~\ref{fig:lowreview}.}
\label{tab:review}
\small
\begin{tabular}{@{}llrrrrrrrr@{}}
\toprule
\textbf{Model} & \textbf{Fam.} & \textbf{Val PR} & \textbf{ROC} &
\textbf{PR-AUC} & \textbf{lift} & \textbf{bal.\ acc} & \textbf{F1} &
\textbf{acc} & \textbf{Brier} \\
\midrule
LogisticRegression \emph{(r0)}   & A & 0.320 & 0.659 & 0.291 & 2.63$\times$ & 0.636 & 0.345 & 0.847 & 0.090 \\
LogisticRegression ($\ell_2$) \emph{(r1)} & A & 0.320 & 0.659 & 0.291 & 2.63$\times$ & 0.636 & 0.342 & 0.843 & 0.090 \\
DecisionTree \emph{(r0)}         & B & 0.389 & 0.722 & 0.396 & 3.59$\times$ & 0.677 & 0.461 & 0.899 & 0.080 \\
RandomForest \emph{(r1)}         & B & 0.403 & 0.744 & 0.405 & 3.67$\times$ & 0.680 & 0.470 & 0.901 & 0.078 \\
HistGradientBoosting \emph{(r2)} & B & 0.407 & 0.749 & 0.415 & 3.76$\times$ & 0.680 & 0.471 & 0.902 & 0.078 \\
VotingClassifier (A+B)           & C & 0.404 & 0.747 & 0.407 & 3.69$\times$ & 0.683 & 0.448 & 0.885 & 0.079 \\
\best{StackingClassifier (A+B)}  & C & \best{0.409} & 0.749 & 0.414 & 3.75$\times$ & 0.681 & 0.470 & 0.901 & 0.078 \\
\bottomrule
\end{tabular}
\end{table}

The selected model reaches \textbf{ROC-AUC 0.749}, \textbf{balanced accuracy
0.681} and \textbf{PR-AUC 0.414 --- 3.75$\times$ the no-skill floor}. Its raw
accuracy is 0.901 against a majority-class rule scoring 0.890, which is exactly
why accuracy is printed but never led with. The tree family beats the linear one
by 1.43$\times$ on PR-AUC --- but a \emph{single depth-limited decision tree}
already reaches 0.396 of the eventual 0.414, so 96\% of the achievable signal
comes from one interpretable tree. And the model does not memorise buyers: on the
12{,}617 test orders from customers never seen in training it scores PR-AUC
\textbf{0.415} against 0.414 overall.

\paragraph{Class weighting: the ablation Phase 1 promised.}
On the best tree model, weighting the loss moves ROC-AUC from 0.749 to 0.747 and
PR-AUC from 0.415 to 0.411 --- no gain --- while the Brier score degrades from
\textbf{0.0778 to 0.1331}, a 71\% deterioration, and mean predicted risk inflates
from 0.119 to \textbf{0.337} against an actual rate of 0.110. The weighted model
overstates every customer's risk threefold (Figure~\ref{fig:calibration}). We
therefore fit on the observed prior and move the threshold instead. Tuned on
validation by maximising F1 it lands at 0.19 rather than 0.50, which on the test
window raises F1 from 0.414 to \textbf{0.470} and recall from 0.297 to 0.398 at a
precision cost of 0.686 to 0.575, flagging 7.6\% of delivered orders instead of
4.8\%. Calibration holds where the volume is: in the two bins containing
11{,}892 of the 12{,}855 test orders the gap between predicted risk and observed
rate is under 0.013.

\subsection{What the cross-validation protocol was telling us}
\label{sec:results-protocol}

\begin{table}[htbp]
\centering
\caption{One model, three protocols. Raw average precision disagrees; the same
numbers divided by each period's no-skill floor do not.
Figure~\ref{fig:protocol}.}
\label{tab:protocol}
\small
\begin{tabular}{@{}lrrrr@{}}
\toprule
\textbf{Protocol} & \textbf{PR-AUC} & \textbf{fold spread} & \textbf{no-skill floor} & \textbf{lift} \\
\midrule
\code{TimeSeriesSplit}, expanding window & 0.562 & $\pm$0.091 & 0.158 & 3.55$\times$ \\
\code{StratifiedGroupKFold}, shuffled    & 0.580 & $\pm$0.006 & 0.158 & 3.66$\times$ \\
Forward test window --- the truth        & 0.414 & ---        & 0.110 & 3.75$\times$ \\
\bottomrule
\end{tabular}
\end{table}

Both cross-validated estimates sit about 0.16 above the forward result, which
looks like catastrophic optimism --- and almost all of it is mechanical. Average
precision's no-skill baseline \emph{is} the positive rate
\citep{saito2015prplot}, and the low-review rate falls from 15.8\% inside the
training period to 11.0\% in the test window; normalised by that floor the three
agree to within 0.2 of a multiple. \textbf{Raw PR-AUC is not comparable across
periods with different base rates; its lift over the floor is} --- and quoting
the raw number alone would have supported either a claim of severe overfitting or
a claim of none. The shuffled protocol adds a second warning: its fold-to-fold
spread is $\pm$0.006 against the temporal protocol's $\pm$0.091, an interval
fifteen times narrower on an estimate no closer to the truth.

\subsection{Does the third family earn its place, and where do the models fail?}
\label{sec:results-ensemble}

Family C is selected on both problems and the honest answer differs. On
\textbf{lead time} it earns its place: the Voting blend improves on the better of
its members by 0.24 days of MAE (5.9\%) and is the only model in
Table~\ref{tab:reg} with a positive $R^2$, for the extrapolation reason above. On
\textbf{review risk} it does not: Stacking wins the rank vote by a nose, then
matches the boosted trees to three decimals on ROC-AUC and comes in 0.001
\emph{below} them on PR-AUC, so a deployment could ship the single boosted model
with no measurable loss. That is worth saying because the ensemble was
\emph{flattered} here: stacking builds meta-features with
\code{cross\_val\_predict}, which needs the inner folds to partition the training
set, and an expanding window partitions nothing --- so the inner CV must be
contiguous \code{KFold} and the meta-learner sees predictions partly made by
models trained on later orders (Appendix~\ref{app:hyperparameters}). Being
flattered and still failing to separate is a safer conclusion than the reverse.

\label{sec:results-segments}
Across geography the review model is remarkably stable: ROC-AUC runs 0.744 to
0.778 across the five Brazilian regions, with lift over each region's own floor
between 3.6$\times$ and 5.0$\times$. The regression is not
(Figure~\ref{fig:segments}): MAE runs from 3.42 days in the Southeast, where most
orders are, to \textbf{6.84 days in the North}, and the bias is positive
everywhere --- $+1.15$ days in the Southeast rising to $+4.77$ in the North. The
model \emph{over-predicts lead time everywhere}, most severely where deliveries
are longest, which is exactly what a model fitted on a slower past does when
fulfilment has improved underneath it, and the strongest argument here for
refitting on a rolling basis.
